\documentclass[
  doc,
  floatsintext,
  longtable,
  a4paper,
  nolmodern,
  notxfonts,
  notimes,
  12pt,
  colorlinks=true,linkcolor=blue,citecolor=blue,urlcolor=blue]{apa7}

\usepackage{amsmath}
\usepackage{amssymb}

\geometry{inner=1in, outer=1in}
\fancyhfoffset[LE,RO]{0cm}


\usepackage[bidi=default]{babel}
\babelprovide[main,import]{spanish}


% get rid of language-specific shorthands (see #6817):
\let\LanguageShortHands\languageshorthands
\def\languageshorthands#1{}

\RequirePackage{longtable}
\RequirePackage{threeparttablex}

\makeatletter
\renewcommand{\paragraph}{\@startsection{paragraph}{4}{\parindent}%
	{0\baselineskip \@plus 0.2ex \@minus 0.2ex}%
	{-.5em}%
	{\normalfont\normalsize\bfseries\typesectitle}}

\renewcommand{\subparagraph}[1]{\@startsection{subparagraph}{5}{0.5em}%
	{0\baselineskip \@plus 0.2ex \@minus 0.2ex}%
	{-\z@\relax}%
	{\normalfont\normalsize\bfseries\itshape\hspace{\parindent}{#1}\textit{\addperi}}{\relax}}
\makeatother




\usepackage{longtable, booktabs, multirow, multicol, colortbl, hhline, caption, array, float, xpatch}
\usepackage{subcaption}


\renewcommand\thesubfigure{\Alph{subfigure}}
\setcounter{topnumber}{2}
\setcounter{bottomnumber}{2}
\setcounter{totalnumber}{4}
\renewcommand{\topfraction}{0.85}
\renewcommand{\bottomfraction}{0.85}
\renewcommand{\textfraction}{0.15}
\renewcommand{\floatpagefraction}{0.7}

\usepackage{tcolorbox}
\tcbuselibrary{listings,theorems, breakable, skins}
\usepackage{fontawesome5}

\definecolor{quarto-callout-color}{HTML}{909090}
\definecolor{quarto-callout-note-color}{HTML}{0758E5}
\definecolor{quarto-callout-important-color}{HTML}{CC1914}
\definecolor{quarto-callout-warning-color}{HTML}{EB9113}
\definecolor{quarto-callout-tip-color}{HTML}{00A047}
\definecolor{quarto-callout-caution-color}{HTML}{FC5300}
\definecolor{quarto-callout-color-frame}{HTML}{ACACAC}
\definecolor{quarto-callout-note-color-frame}{HTML}{4582EC}
\definecolor{quarto-callout-important-color-frame}{HTML}{D9534F}
\definecolor{quarto-callout-warning-color-frame}{HTML}{F0AD4E}
\definecolor{quarto-callout-tip-color-frame}{HTML}{02B875}
\definecolor{quarto-callout-caution-color-frame}{HTML}{FD7E14}

%\newlength\Oldarrayrulewidth
%\newlength\Oldtabcolsep


\usepackage{hyperref}



\usepackage{color}
\usepackage{fancyvrb}
\newcommand{\VerbBar}{|}
\newcommand{\VERB}{\Verb[commandchars=\\\{\}]}
\DefineVerbatimEnvironment{Highlighting}{Verbatim}{commandchars=\\\{\}}
% Add ',fontsize=\small' for more characters per line
\usepackage{framed}
\definecolor{shadecolor}{RGB}{241,243,245}
\newenvironment{Shaded}{\begin{snugshade}}{\end{snugshade}}
\newcommand{\AlertTok}[1]{\textcolor[rgb]{0.68,0.00,0.00}{#1}}
\newcommand{\AnnotationTok}[1]{\textcolor[rgb]{0.37,0.37,0.37}{#1}}
\newcommand{\AttributeTok}[1]{\textcolor[rgb]{0.40,0.45,0.13}{#1}}
\newcommand{\BaseNTok}[1]{\textcolor[rgb]{0.68,0.00,0.00}{#1}}
\newcommand{\BuiltInTok}[1]{\textcolor[rgb]{0.00,0.23,0.31}{#1}}
\newcommand{\CharTok}[1]{\textcolor[rgb]{0.13,0.47,0.30}{#1}}
\newcommand{\CommentTok}[1]{\textcolor[rgb]{0.37,0.37,0.37}{#1}}
\newcommand{\CommentVarTok}[1]{\textcolor[rgb]{0.37,0.37,0.37}{\textit{#1}}}
\newcommand{\ConstantTok}[1]{\textcolor[rgb]{0.56,0.35,0.01}{#1}}
\newcommand{\ControlFlowTok}[1]{\textcolor[rgb]{0.00,0.23,0.31}{\textbf{#1}}}
\newcommand{\DataTypeTok}[1]{\textcolor[rgb]{0.68,0.00,0.00}{#1}}
\newcommand{\DecValTok}[1]{\textcolor[rgb]{0.68,0.00,0.00}{#1}}
\newcommand{\DocumentationTok}[1]{\textcolor[rgb]{0.37,0.37,0.37}{\textit{#1}}}
\newcommand{\ErrorTok}[1]{\textcolor[rgb]{0.68,0.00,0.00}{#1}}
\newcommand{\ExtensionTok}[1]{\textcolor[rgb]{0.00,0.23,0.31}{#1}}
\newcommand{\FloatTok}[1]{\textcolor[rgb]{0.68,0.00,0.00}{#1}}
\newcommand{\FunctionTok}[1]{\textcolor[rgb]{0.28,0.35,0.67}{#1}}
\newcommand{\ImportTok}[1]{\textcolor[rgb]{0.00,0.46,0.62}{#1}}
\newcommand{\InformationTok}[1]{\textcolor[rgb]{0.37,0.37,0.37}{#1}}
\newcommand{\KeywordTok}[1]{\textcolor[rgb]{0.00,0.23,0.31}{\textbf{#1}}}
\newcommand{\NormalTok}[1]{\textcolor[rgb]{0.00,0.23,0.31}{#1}}
\newcommand{\OperatorTok}[1]{\textcolor[rgb]{0.37,0.37,0.37}{#1}}
\newcommand{\OtherTok}[1]{\textcolor[rgb]{0.00,0.23,0.31}{#1}}
\newcommand{\PreprocessorTok}[1]{\textcolor[rgb]{0.68,0.00,0.00}{#1}}
\newcommand{\RegionMarkerTok}[1]{\textcolor[rgb]{0.00,0.23,0.31}{#1}}
\newcommand{\SpecialCharTok}[1]{\textcolor[rgb]{0.37,0.37,0.37}{#1}}
\newcommand{\SpecialStringTok}[1]{\textcolor[rgb]{0.13,0.47,0.30}{#1}}
\newcommand{\StringTok}[1]{\textcolor[rgb]{0.13,0.47,0.30}{#1}}
\newcommand{\VariableTok}[1]{\textcolor[rgb]{0.07,0.07,0.07}{#1}}
\newcommand{\VerbatimStringTok}[1]{\textcolor[rgb]{0.13,0.47,0.30}{#1}}
\newcommand{\WarningTok}[1]{\textcolor[rgb]{0.37,0.37,0.37}{\textit{#1}}}

\providecommand{\tightlist}{%
  \setlength{\itemsep}{0pt}\setlength{\parskip}{0pt}}
\usepackage{longtable,booktabs,array}
\newcounter{none} % for unnumbered tables
\usepackage{calc} % for calculating minipage widths
% Correct order of tables after \paragraph or \subparagraph
\usepackage{etoolbox}
\makeatletter
\patchcmd\longtable{\par}{\if@noskipsec\mbox{}\fi\par}{}{}
\makeatother
% Allow footnotes in longtable head/foot
\IfFileExists{footnotehyper.sty}{\usepackage{footnotehyper}}{\usepackage{footnote}}
\makesavenoteenv{longtable}

\usepackage{graphicx}
\makeatletter
\newsavebox\pandoc@box
\newcommand*\pandocbounded[1]{% scales image to fit in text height/width
  \sbox\pandoc@box{#1}%
  \Gscale@div\@tempa{\textheight}{\dimexpr\ht\pandoc@box+\dp\pandoc@box\relax}%
  \Gscale@div\@tempb{\linewidth}{\wd\pandoc@box}%
  \ifdim\@tempb\p@<\@tempa\p@\let\@tempa\@tempb\fi% select the smaller of both
  \ifdim\@tempa\p@<\p@\scalebox{\@tempa}{\usebox\pandoc@box}%
  \else\usebox{\pandoc@box}%
  \fi%
}
% Set default figure placement to htbp
\def\fps@figure{htbp}
\makeatother







\usepackage{newtx}

\defaultfontfeatures{Scale=MatchLowercase}
\defaultfontfeatures[\rmfamily]{Ligatures=TeX,Scale=1}





\title{Variables globales}


\shorttitle{VARIABLES GLOBALES}


\usepackage{etoolbox}



\ccoppy{\textcopyright~2026}



\author{Edison Achalma}



\affiliation{
{Escuela Profesional de Economía, Universidad Nacional de San Cristóbal
de Huamanga}}




\leftheader{Achalma}

\date{2026-10-20}


\abstract{Este abstract será actualizado una vez que se complete el
contenido final del artículo. }



\authornote{\par{\addORCIDlink{Edison Achalma}{0000-0001-6996-3364}} 
\par{ }
\par{   El autor no tiene conflictos de interés que revelar.    Los
roles de autor se clasificaron utilizando la taxonomía de roles de
colaborador (CRediT; https://credit.niso.org/) de la siguiente
manera:  Edison Achalma:   conceptualización, metodología, análisis
formal, investigación, recursos, curación de
datos, redacción, visualización, supervisión, administración del
proyecto}
\par{La correspondencia relativa a este artículo debe dirigirse a Edison
Achalma, Escuela Profesional de Economía, Universidad Nacional de San
Cristóbal de Huamanga, Portal Independencia N°
57, Ayacucho, AYA 5001, Perú, Email: \href{mailto:elmer.achalma.09@unsch.edu.pe}{elmer.achalma.09@unsch.edu.pe}}
}

\makeatletter
\let\endoldlt\endlongtable
\def\endlongtable{
\hline
\endoldlt
}
\makeatother

\urlstyle{same}



\hypersetup{
  pdftitle={Variables globales},
  pdfauthor={Autor}
}
\makeatletter
\@ifpackageloaded{caption}{}{\usepackage{caption}}
\AtBeginDocument{%
\ifdefined\contentsname
  \renewcommand*\contentsname{Tabla de contenidos}
\else
  \newcommand\contentsname{Tabla de contenidos}
\fi
\ifdefined\listfigurename
  \renewcommand*\listfigurename{Lista de Figuras}
\else
  \newcommand\listfigurename{Lista de Figuras}
\fi
\ifdefined\listtablename
  \renewcommand*\listtablename{Lista de Tablas}
\else
  \newcommand\listtablename{Lista de Tablas}
\fi
\ifdefined\figurename
  \renewcommand*\figurename{Figura}
\else
  \newcommand\figurename{Figura}
\fi
\ifdefined\tablename
  \renewcommand*\tablename{Tabla}
\else
  \newcommand\tablename{Tabla}
\fi
}
\@ifpackageloaded{float}{}{\usepackage{float}}
\floatstyle{ruled}
\@ifundefined{c@chapter}{\newfloat{codelisting}{h}{lop}}{\newfloat{codelisting}{h}{lop}[chapter]}
\floatname{codelisting}{Listado}
\newcommand*\listoflistings{\listof{codelisting}{Lista de Listados}}
\makeatother
\makeatletter
\makeatother
\makeatletter
\@ifpackageloaded{caption}{}{\usepackage{caption}}
\@ifpackageloaded{subcaption}{}{\usepackage{subcaption}}
\makeatother
\makeatletter
\@ifpackageloaded{fontawesome5}{}{\usepackage{fontawesome5}}
\makeatother

% From https://tex.stackexchange.com/a/645996/211326
%%% apa7 doesn't want to add appendix section titles in the toc
%%% let's make it do it
\makeatletter
\xpatchcmd{\appendix}
  {\par}
  {\addcontentsline{toc}{section}{\@currentlabelname}\par}
  {}{}
\makeatother

%% Disable longtable counter
%% https://tex.stackexchange.com/a/248395/211326

\usepackage{etoolbox}

\makeatletter
\patchcmd{\LT@caption}
  {\bgroup}
  {\bgroup\global\LTpatch@captiontrue}
  {}{}
\patchcmd{\longtable}
  {\par}
  {\par\global\LTpatch@captionfalse}
  {}{}
\apptocmd{\endlongtable}
  {\ifLTpatch@caption\else\addtocounter{table}{-1}\fi}
  {}{}
\newif\ifLTpatch@caption
\makeatother

\begin{document}

\maketitle


\hypertarget{toc}{}
\tableofcontents
\newpage
\section[Introduction]{Variables globales}

\setcounter{secnumdepth}{3}

\setlength\LTleft{0pt}




\subsection{🎯 Motivación}\label{motivaciuxf3n}

Cuando trabajas con funciones, necesitas entender cómo Python maneja las
variables en diferentes contextos. Las variables globales son accesibles
desde cualquier parte de tu programa, pero modificarlas dentro de
funciones requiere cuidado. Dominar el alcance (scope) de las variables
es esencial para escribir código predecible y evitar errores difíciles
de depurar.

\subsection{💡 Intuición}\label{intuiciuxf3n}

Imagina que tienes una pizarra en la entrada de tu casa (variable
global) donde toda la familia puede escribir mensajes. Todos pueden leer
lo que está en la pizarra. Pero si alguien escribe un mensaje en su
propia libreta personal (variable local) con el mismo título, no afecta
lo que está en la pizarra. Python funciona igual: una variable dentro de
una función es como una libreta privada, separada de la pizarra global.

\subsection{🧠 Modelo mental}\label{modelo-mental}

Python sigue las reglas \textbf{LEGB} para la resolución de variables:

\begin{enumerate}
\def\labelenumi{\arabic{enumi}.}
\tightlist
\item
  \textbf{L}ocal → dentro de la función actual
\item
  \textbf{E}nclosing → funciones externas (anidadas)
\item
  \textbf{G}lobal → variables definidas fuera de funciones
\item
  \textbf{B}uilt-in → funciones y variables incorporadas de Python
\end{enumerate}

Cuando usas una variable dentro de una función, Python busca en ese
orden: primero local, luego global si no encuentra local.

\subsection{📐 Desarrollo}\label{desarrollo}

\subsubsection{1. Variables globales
(acceso)}\label{variables-globales-acceso}

Las variables creadas fuera de funciones son globales y accesibles
dentro de funciones:

\begin{Shaded}
\begin{Highlighting}[]
\NormalTok{x }\OperatorTok{=} \StringTok{"awesome"}  \CommentTok{\# Variable global}

\KeywordTok{def}\NormalTok{ myfunc():}
    \BuiltInTok{print}\NormalTok{(}\StringTok{"Python is "} \OperatorTok{+}\NormalTok{ x)  }\CommentTok{\# Accede a la variable global}

\NormalTok{myfunc()  }\CommentTok{\# Salida: Python is awesome}
\end{Highlighting}
\end{Shaded}

\subsubsection{2. Variable local con mismo
nombre}\label{variable-local-con-mismo-nombre}

Si defines una variable con el mismo nombre dentro de una función,
\textbf{es una variable local independiente}:

\begin{Shaded}
\begin{Highlighting}[]
\NormalTok{x }\OperatorTok{=} \StringTok{"awesome"}  \CommentTok{\# Global}

\KeywordTok{def}\NormalTok{ myfunc():}
\NormalTok{    x }\OperatorTok{=} \StringTok{"fantastic"}  \CommentTok{\# Local (independiente)}
    \BuiltInTok{print}\NormalTok{(}\StringTok{"Python is "} \OperatorTok{+}\NormalTok{ x)  }\CommentTok{\# Usa la local}

\NormalTok{myfunc()  }\CommentTok{\# Salida: Python is fantastic}
\BuiltInTok{print}\NormalTok{(}\StringTok{"Python is "} \OperatorTok{+}\NormalTok{ x)  }\CommentTok{\# Salida: Python is awesome (global)}
\end{Highlighting}
\end{Shaded}

\subsubsection{\texorpdfstring{3. La palabra clave
\texttt{global}}{3. La palabra clave global}}\label{la-palabra-clave-global}

Para \textbf{modificar} una variable global dentro de una función, usa
\texttt{global}:

\begin{Shaded}
\begin{Highlighting}[]
\NormalTok{x }\OperatorTok{=} \StringTok{"awesome"}  \CommentTok{\# Global}

\KeywordTok{def}\NormalTok{ myfunc():}
    \KeywordTok{global}\NormalTok{ x  }\CommentTok{\# Declara que vas a usar/modificar la global}
\NormalTok{    x }\OperatorTok{=} \StringTok{"fantastic"}  \CommentTok{\# Modifica la global}

\NormalTok{myfunc()}
\BuiltInTok{print}\NormalTok{(}\StringTok{"Python is "} \OperatorTok{+}\NormalTok{ x)  }\CommentTok{\# Salida: Python is fantastic}
\end{Highlighting}
\end{Shaded}

\subsubsection{4. Crear variable global dentro de
función}\label{crear-variable-global-dentro-de-funciuxf3n}

Puedes crear una nueva variable global desde dentro de una función:

\begin{Shaded}
\begin{Highlighting}[]
\KeywordTok{def}\NormalTok{ myfunc():}
    \KeywordTok{global}\NormalTok{ x}
\NormalTok{    x }\OperatorTok{=} \StringTok{"nueva variable"}  \CommentTok{\# Crea una variable global}

\NormalTok{myfunc()}
\BuiltInTok{print}\NormalTok{(x)  }\CommentTok{\# Salida: nueva variable}
\end{Highlighting}
\end{Shaded}

\subsubsection{5. Modificar variables globales vs
locales}\label{modificar-variables-globales-vs-locales}

\begin{Shaded}
\begin{Highlighting}[]
\NormalTok{contador }\OperatorTok{=} \DecValTok{0}  \CommentTok{\# Global}

\KeywordTok{def}\NormalTok{ incrementar():}
    \KeywordTok{global}\NormalTok{ contador  }\CommentTok{\# Necesario para modificar}
\NormalTok{    contador }\OperatorTok{+=} \DecValTok{1}

\KeywordTok{def}\NormalTok{ mostrar():}
    \BuiltInTok{print}\NormalTok{(contador)  }\CommentTok{\# Solo lectura, no necesita global}

\NormalTok{incrementar()}
\NormalTok{incrementar()}
\NormalTok{mostrar()  }\CommentTok{\# Salida: 2}
\end{Highlighting}
\end{Shaded}

\subsection{🔬 Ejemplos}\label{ejemplos}

\textbf{Mínimo ---} Acceso a global:

\begin{Shaded}
\begin{Highlighting}[]
\NormalTok{nombre }\OperatorTok{=} \StringTok{"Ana"}  \CommentTok{\# Global}

\KeywordTok{def}\NormalTok{ saludar():}
    \BuiltInTok{print}\NormalTok{(}\StringTok{"Hola"}\NormalTok{, nombre)  }\CommentTok{\# Accede a la global}

\NormalTok{saludar()  }\CommentTok{\# Hola Ana}
\end{Highlighting}
\end{Shaded}

\textbf{Progresivo ---} Local vs Global:

\begin{Shaded}
\begin{Highlighting}[]
\NormalTok{mensaje }\OperatorTok{=} \StringTok{"Global"}

\KeywordTok{def}\NormalTok{ funcion1():}
\NormalTok{    mensaje }\OperatorTok{=} \StringTok{"Local"}
    \BuiltInTok{print}\NormalTok{(}\StringTok{"Dentro:"}\NormalTok{, mensaje)  }\CommentTok{\# Local}

\KeywordTok{def}\NormalTok{ funcion2():}
    \BuiltInTok{print}\NormalTok{(}\StringTok{"Fuera:"}\NormalTok{, mensaje)  }\CommentTok{\# Global}

\NormalTok{funcion1()  }\CommentTok{\# Dentro: Local}
\NormalTok{funcion2()  }\CommentTok{\# Fuera: Global}
\BuiltInTok{print}\NormalTok{(}\StringTok{"Principal:"}\NormalTok{, mensaje)  }\CommentTok{\# Principal: Global}
\end{Highlighting}
\end{Shaded}

\textbf{Real ---} Contador de visitas:

\begin{Shaded}
\begin{Highlighting}[]
\NormalTok{visitas }\OperatorTok{=} \DecValTok{0}  \CommentTok{\# Global}

\KeywordTok{def}\NormalTok{ visitar\_pagina():}
    \KeywordTok{global}\NormalTok{ visitas}
\NormalTok{    visitas }\OperatorTok{+=} \DecValTok{1}
    \BuiltInTok{print}\NormalTok{(}\SpecialStringTok{f"Visitante \#}\SpecialCharTok{\{}\NormalTok{visitas}\SpecialCharTok{\}}\SpecialStringTok{"}\NormalTok{)}

\KeywordTok{def}\NormalTok{ mostrar\_estadisticas():}
    \BuiltInTok{print}\NormalTok{(}\SpecialStringTok{f"Total de visitas: }\SpecialCharTok{\{}\NormalTok{visitas}\SpecialCharTok{\}}\SpecialStringTok{"}\NormalTok{)}

\CommentTok{\# Simulación de visitas}
\NormalTok{visitar\_pagina()  }\CommentTok{\# Visitante \#1}
\NormalTok{visitar\_pagina()  }\CommentTok{\# Visitante \#2}
\NormalTok{visitar\_pagina()  }\CommentTok{\# Visitante \#3}
\NormalTok{mostrar\_estadisticas()  }\CommentTok{\# Total de visitas: 3}
\end{Highlighting}
\end{Shaded}

\subsection{⚠️ Errores comunes}\label{errores-comunes}

{\def\LTcaptype{none} % do not increment counter
\begin{longtable}[]{@{}
  >{\raggedright\arraybackslash}p{(\linewidth - 4\tabcolsep) * \real{0.2920}}
  >{\raggedright\arraybackslash}p{(\linewidth - 4\tabcolsep) * \real{0.3451}}
  >{\raggedright\arraybackslash}p{(\linewidth - 4\tabcolsep) * \real{0.3628}}@{}}
\toprule\noalign{}
\begin{minipage}[b]{\linewidth}\raggedright
Error
\end{minipage} & \begin{minipage}[b]{\linewidth}\raggedright
Síntoma
\end{minipage} & \begin{minipage}[b]{\linewidth}\raggedright
Solución
\end{minipage} \\
\midrule\noalign{}
\endhead
\bottomrule\noalign{}
\endlastfoot
\textbf{Modificar global sin \texttt{global}} & Se crea una local en
lugar de modificar & Añade \texttt{global\ variable} \\
\textbf{Usar variable sin inicializar} & \texttt{UnboundLocalError} &
Inicializa la variable antes de usarla \\
\textbf{Confundir scope} & La variable no tiene el valor esperado & Usa
\texttt{global} para modificar, no para leer \\
\textbf{Crear muchas globales} & Código difícil de seguir & Usa
parámetros en lugar de globales \\
\end{longtable}
}

\subsection{✅ Buenas prácticas}\label{buenas-pruxe1cticas}

\begin{enumerate}
\def\labelenumi{\arabic{enumi}.}
\item
  \textbf{Evita usar variables globales a menos que sea necesario}: pasa
  valores como argumentos a las funciones.
\item
  \textbf{Usa globales solo para configuraciones o contadores}:

\begin{Shaded}
\begin{Highlighting}[]
\CommentTok{\# ✅ Aceptable}
\NormalTok{CONFIG }\OperatorTok{=}\NormalTok{ \{}\StringTok{"debug"}\NormalTok{: }\VariableTok{True}\NormalTok{, }\StringTok{"timeout"}\NormalTok{: }\DecValTok{30}\NormalTok{\}}

\CommentTok{\# ✅ Buen uso de global}
\NormalTok{visitas }\OperatorTok{=} \DecValTok{0}
\KeywordTok{def}\NormalTok{ incrementar\_visitas():}
    \KeywordTok{global}\NormalTok{ visitas}
\NormalTok{    visitas }\OperatorTok{+=} \DecValTok{1}
\end{Highlighting}
\end{Shaded}
\item
  \textbf{Nombra las globales con mayúsculas} (por convención):

\begin{Shaded}
\begin{Highlighting}[]
\NormalTok{MAX\_INTENTOS }\OperatorTok{=} \DecValTok{3}  \CommentTok{\# Constante global}
\NormalTok{BASE\_URL }\OperatorTok{=} \StringTok{"https://api.ejemplo.com"}  \CommentTok{\# Configuración}
\end{Highlighting}
\end{Shaded}
\item
  \textbf{Si necesitas modificar una variable, considera pasar por
  parámetros}:

\begin{Shaded}
\begin{Highlighting}[]
\CommentTok{\# ❌ Uso global innecesario}
\NormalTok{saldo }\OperatorTok{=} \DecValTok{100}
\KeywordTok{def}\NormalTok{ sumar(monto):}
    \KeywordTok{global}\NormalTok{ saldo}
\NormalTok{    saldo }\OperatorTok{+=}\NormalTok{ monto}

\CommentTok{\# ✅ Mejor usar parámetros}
\KeywordTok{def}\NormalTok{ sumar(saldo, monto):}
    \ControlFlowTok{return}\NormalTok{ saldo }\OperatorTok{+}\NormalTok{ monto}
\end{Highlighting}
\end{Shaded}
\item
  \textbf{Documenta las variables globales} para que sea claro qué
  hacen.
\end{enumerate}

\subsection{📝 Ejercicio}\label{ejercicio}

\textbf{Básico} --- Análisis de alcance

Analiza el siguiente código y determina qué imprime cada
\texttt{print()}:

\begin{Shaded}
\begin{Highlighting}[]
\NormalTok{mensaje }\OperatorTok{=} \StringTok{"Hola mundo"}

\KeywordTok{def}\NormalTok{ funcion\_a():}
\NormalTok{    mensaje }\OperatorTok{=} \StringTok{"Hola desde A"}
    \BuiltInTok{print}\NormalTok{(mensaje)  }\CommentTok{\# 1}

\KeywordTok{def}\NormalTok{ funcion\_b():}
    \KeywordTok{global}\NormalTok{ mensaje}
\NormalTok{    mensaje }\OperatorTok{=} \StringTok{"Modificado globalmente"}
    \BuiltInTok{print}\NormalTok{(mensaje)  }\CommentTok{\# 2}

\KeywordTok{def}\NormalTok{ funcion\_c():}
    \BuiltInTok{print}\NormalTok{(mensaje)  }\CommentTok{\# 3}

\NormalTok{funcion\_a()}
\BuiltInTok{print}\NormalTok{(mensaje)  }\CommentTok{\# 4}
\NormalTok{funcion\_b()}
\BuiltInTok{print}\NormalTok{(mensaje)  }\CommentTok{\# 5}
\NormalTok{funcion\_c()}
\end{Highlighting}
\end{Shaded}

\begin{center}\rule{0.5\linewidth}{0.5pt}\end{center}

\textbf{Solución (no mires hasta haberlo intentado):}

\begin{Shaded}
\begin{Highlighting}[]
\NormalTok{mensaje }\OperatorTok{=} \StringTok{"Hola mundo"}

\KeywordTok{def}\NormalTok{ funcion\_a():}
\NormalTok{    mensaje }\OperatorTok{=} \StringTok{"Hola desde A"}  \CommentTok{\# Variable local}
    \BuiltInTok{print}\NormalTok{(mensaje)  }\CommentTok{\# 1. Imprime: Hola desde A}

\KeywordTok{def}\NormalTok{ funcion\_b():}
    \KeywordTok{global}\NormalTok{ mensaje  }\CommentTok{\# Acceso a la global}
\NormalTok{    mensaje }\OperatorTok{=} \StringTok{"Modificado globalmente"}  \CommentTok{\# Cambia la global}
    \BuiltInTok{print}\NormalTok{(mensaje)  }\CommentTok{\# 2. Imprime: Modificado globalmente}

\KeywordTok{def}\NormalTok{ funcion\_c():}
    \BuiltInTok{print}\NormalTok{(mensaje)  }\CommentTok{\# 3. Imprime la global (que es: Modificado globalmente)}

\NormalTok{funcion\_a()}
\CommentTok{\# 4. Imprime: Hola mundo (la global no cambió)}

\NormalTok{funcion\_b()}
\CommentTok{\# 5. Imprime: Modificado globalmente (la global cambió)}

\NormalTok{funcion\_c()}
\CommentTok{\# 6. Imprime: Modificado globalmente (la global sigue igual)}
\end{Highlighting}
\end{Shaded}

\textbf{Salida esperada:}

\begin{verbatim}
Hola desde A
Hola mundo
Modificado globalmente
Modificado globalmente
Modificado globalmente
\end{verbatim}

\subsection{🗂️ Síntesis}\label{suxedntesis}

\begin{itemize}
\tightlist
\item
  \textbf{Global}: variable definida fuera de funciones, accesible en
  todo el programa.
\item
  \textbf{Local}: variable definida dentro de una función, solo
  accesible allí.
\item
  \textbf{Leer global} no requiere \texttt{global}; solo usas el nombre.
\item
  \textbf{Modificar global} requiere \texttt{global\ variable}.
\item
  \textbf{Crear global} dentro de función: usa \texttt{global} para
  crearla.
\item
  \textbf{\texttt{global}} declara que la variable pertenece al scope
  global.
\item
  \textbf{Buena práctica}: minimiza el uso de variables globales.
\end{itemize}

\subsection{🔗 Qué sigue}\label{quuxe9-sigue}

\begin{itemize}
\tightlist
\item
  \textbf{Relacionado:} funciones (alcance de variables, parámetros y
  argumentos)
\item
  \textbf{Siguiente tema:} Funciones --- aprenderás a crear y usar
  funciones, y a manejar correctamente el alcance de variables en
  diferentes contextos.
\end{itemize}

\section{Publicaciones Similares}\label{publicaciones-similares}

Si te interesó este artículo, te recomendamos que explores otros blogs y
recursos relacionados que pueden ampliar tus conocimientos. Aquí te dejo
algunas sugerencias:

\begin{enumerate}
\def\labelenumi{\arabic{enumi}.}
\tightlist
\item
  \href{https://numerus-scriptum.netlify.app/python/2020-06-19-instalacion-de-anaconda/index.pdf}{\faIcon{file-pdf}}
  \href{https://numerus-scriptum.netlify.app/python/2020-06-19-instalacion-de-anaconda}{Instalacion
  De Anaconda}
\item
  \href{https://numerus-scriptum.netlify.app/python/2020-06-20-configurar-entorno-virtual-python-anaconda/index.pdf}{\faIcon{file-pdf}}
  \href{https://numerus-scriptum.netlify.app/python/2020-06-20-configurar-entorno-virtual-python-anaconda}{Configurar
  Entorno Virtual Python Anaconda}
\item
  \href{https://numerus-scriptum.netlify.app/python/2021-04-17-01-introducion-a-la-programacion-con-python/index.pdf}{\faIcon{file-pdf}}
  \href{https://numerus-scriptum.netlify.app/python/2021-04-17-01-introducion-a-la-programacion-con-python}{01
  Introducion A La Programacion Con Python}
\item
  \href{https://numerus-scriptum.netlify.app/python/2021-05-31-02-variables-expresiones-y-statements-con-python/index.pdf}{\faIcon{file-pdf}}
  \href{https://numerus-scriptum.netlify.app/python/2021-05-31-02-variables-expresiones-y-statements-con-python}{02
  Variables Expresiones Y Statements Con Python}
\item
  \href{https://numerus-scriptum.netlify.app/python/2021-06-07-03-objetos-de-python/index.pdf}{\faIcon{file-pdf}}
  \href{https://numerus-scriptum.netlify.app/python/2021-06-07-03-objetos-de-python}{03
  Objetos De Python}
\item
  \href{https://numerus-scriptum.netlify.app/python/2021-06-14-04-ejecucion-condicional-con-python/index.pdf}{\faIcon{file-pdf}}
  \href{https://numerus-scriptum.netlify.app/python/2021-06-14-04-ejecucion-condicional-con-python}{04
  Ejecucion Condicional Con Python}
\item
  \href{https://numerus-scriptum.netlify.app/python/2021-06-21-05-iteraciones-con-python/index.pdf}{\faIcon{file-pdf}}
  \href{https://numerus-scriptum.netlify.app/python/2021-06-21-05-iteraciones-con-python}{05
  Iteraciones Con Python}
\item
  \href{https://numerus-scriptum.netlify.app/python/2021-08-16-06-funciones-con-python/index.pdf}{\faIcon{file-pdf}}
  \href{https://numerus-scriptum.netlify.app/python/2021-08-16-06-funciones-con-python}{06
  Funciones Con Python}
\item
  \href{https://numerus-scriptum.netlify.app/python/2021-08-23-07-dataframes-con-python/index.pdf}{\faIcon{file-pdf}}
  \href{https://numerus-scriptum.netlify.app/python/2021-08-23-07-dataframes-con-python}{07
  Dataframes Con Python}
\item
  \href{https://numerus-scriptum.netlify.app/python/2021-11-29-08-prediccion-y-metrica-de-performance-con-python/index.pdf}{\faIcon{file-pdf}}
  \href{https://numerus-scriptum.netlify.app/python/2021-11-29-08-prediccion-y-metrica-de-performance-con-python}{08
  Prediccion Y Metrica De Performance Con Python}
\item
  \href{https://numerus-scriptum.netlify.app/python/2021-12-06-09-metodos-de-machine-learning-para-clasificacion-con-python/index.pdf}{\faIcon{file-pdf}}
  \href{https://numerus-scriptum.netlify.app/python/2021-12-06-09-metodos-de-machine-learning-para-clasificacion-con-python}{09
  Metodos De Machine Learning Para Clasificacion Con Python}
\item
  \href{https://numerus-scriptum.netlify.app/python/2021-12-13-10-metodos-de-machine-learning-para-regresion-con-python/index.pdf}{\faIcon{file-pdf}}
  \href{https://numerus-scriptum.netlify.app/python/2021-12-13-10-metodos-de-machine-learning-para-regresion-con-python}{10
  Metodos De Machine Learning Para Regresion Con Python}
\item
  \href{https://numerus-scriptum.netlify.app/python/2025-05-10-visualizacion-de-datos-con-python/index.pdf}{\faIcon{file-pdf}}
  \href{https://numerus-scriptum.netlify.app/python/2025-05-10-visualizacion-de-datos-con-python}{Visualizacion
  De Datos Con Python}
\end{enumerate}

Esperamos que encuentres estas publicaciones igualmente interesantes y
útiles. ¡Disfruta de la lectura!






\end{document}
